\subsection{Wolfram cellular automata main result}\label{sec:wolfram}

We evaluate the third domain on Wolfram cellular automata \cite{wolfram2002}, one-dimensional discrete dynamical systems in which each cell's binary state is updated by a fixed local rule based on neighboring states. For each seed, we sample 200 training transitions and 50 evaluation transitions, with maximum episode length 16. We train for 200 epochs, longer than the molecular experiments, because the discrete dynamics are more challenging under the same flat world-model architecture. We evaluate 10 seeds per prior at shared weight \(\lambda=0.01\), using the default per-step Laplacian construction from the cell adjacency graph for the spectral prior.

Table~\ref{tab:wolfram_main} and Figure~\ref{fig:wolfram_main} summarize the main Wolfram result. With ten seeds per prior, the spectral prior has mean H=16 error \(1.109 \pm 2.620\) and median \(0.113\), compared with mean \(0.069 \pm 0.038\) and median \(0.052\) for no prior, and mean \(0.057 \pm 0.020\) and median \(0.054\) for the Euclidean prior. The mean comparison is dominated by heavy-tail behavior rather than uniform degradation: the per-seed spectral H=16 values, sorted, are \([0.044, 0.060, 0.076, 0.085, 0.091, 0.135, 0.348, 0.462, 0.854, 8.934]\). Six seeds remain at or below the no-prior baseline range, three fall in the [0.35, 0.85] regime (5-15x baseline but still converged), and one seed diverges to 8.934. Median (0.113) and interquartile range therefore characterize the typical seed more faithfully than the mean. We further note that at shorter horizons \(H \in \{1, 2, 4\}\), the spectral median actually lies \emph{below} the no-prior baseline (Figure~\ref{fig:wolfram_main}), indicating the prior is helpful in the short-horizon regime and only diverges as horizon and the catastrophic minority of seeds compound. The Euclidean prior is more stable on this domain, with mean \(0.057\), median \(0.054\), and maximum \(0.088\), comparable to the no-prior baseline.

At long horizons, the spectral prior at \(\lambda=0.01\) does not improve over the no-prior baseline on Wolfram CA, in contrast to the rMD17 and ISO17 settings (where spectral cuts long-horizon error by 51-57\%). We report this as a negative result at the shared weight. The important feature is not that every spectral run degrades, but that the distribution is fragile: a minority of seeds enters a divergent long-horizon regime and controls the mean. Section~\ref{sec:wolfram_mechanism} shows that this is not a fundamental incompatibility between spectral regularization and cellular automata. With a smaller weight, \(\lambda=0.005\), and per-step Laplacian recomputation, the spectral prior matches the no-prior baseline at H=16 (\(0.062\) versus \(0.066\)). Thus, \(\lambda=0.01\) appears to lie in a transition zone where most seeds remain stable but occasional seeds produce large rollout errors.

We therefore do not interpret the Wolfram result as evidence against the spectral prior in general. Rather, the molecular successes in Sections~\ref{sec:rmd17} and~\ref{sec:iso17}, together with the Wolfram mechanism ablation in Section~\ref{sec:wolfram_mechanism}, suggest a conditional rule: spectral priors help when the Laplacian construction matches the data's graph evolution, and become fragile when the prior is too strong or the graph approximation is stale. Section~\ref{sec:discussion} returns to this interpretation as the central mechanism statement of the paper.
